% !PRINTED 631
\weberpage{631}
If we make the substitution $(r, r^n)$ in it, this becomes
\weq{14}{r^n \equiv g^{-\frac{p-1}{m}} \ (\mathrm{mod}\ \mathfrak p_n),}
or, if we define $n'$ by the congruence
\weq{15}{n n' \equiv 1 \ (\mathrm{mod}\ m),}
then:
\weq{16}{r \equiv g^{-n'\frac{p-1}{m}} \ (\mathrm{mod}\ \mathfrak p_n).}

In this case, therefore, in agreement with the general
 rule, $\mathfrak p_n$ is the greatest common divisor of
\[
p \ \text{and} \ \left(r - g^{-n'\frac{p-1}{m}}\right).
\]

By this determination the individual prime forms $\mathfrak p_n$
are exactly characterised. We thus obtain the theorem:

2. \gesperrt{A prime $p$ which is congruent to 1 with respect
to the modulus $m$ splits in the field $\Omega_m$ into $\varphi(m)$
prime factors $\mathfrak p_n$ of the first degree, distinct
from one another, which one obtains as the greatest
 common divisors of $p$ with the various
numbers $r - g^{-n'\frac{p-1}{m}}$, if $g$ is a primitive
root of $p$ and $n'$ is determined by the congruence
$n n' \equiv 1 \ (\mathrm{mod}\ m)$.}

According to \S.\,141 we can decompose every integer or fractional
number $\omega$, and every functional of the field $\Omega_m$, into prime
 factors in such a way that numerator and denominator contain
 no common prime factor, and this decomposition is,
apart from unit factors, completely determined, so that, if
$\mathfrak p$ is a prime ideal, there exists a perfectly determinate positive
or negative exponent $k$ such that the functional $\omega \mathfrak p^{-k}$ contains
 the prime factor $k$ neither in the numerator nor in the denominator. % [sic] source prints $k$; context requires $\mathfrak p$
We then say,
for the sake of brevity, that $\mathfrak p^k$ is \gesperrt{the power of $\mathfrak p$ contained
in $\omega$}.

If $k=0$, we say that $\mathfrak p$ is \gesperrt{not contained} in $\omega$.

Two numbers (or also functionals) $\omega, \omega'$ of a field, in
which one and the same prime functional $\mathfrak p$ is contained, \gesperrt{we
shall call divisor-related, or related for short}; if, on the other hand,
no prime functional is contained in both $\omega$ and $\omega'$, then
$\omega$ and $\omega'$ are called \gesperrt{coprime, or foreign}. This expression
will be particularly useful to us in what follows, when
%CONT>>
